\section{Discussion}
\label{sec:discussion}

\subsection{Adaptivity, not polarisation, decides accuracy}

The four-detector comparison changes how the single- versus dual-polarisation
question should be posed. When both detectors adapt to each scene, the
dual-polarisation SVM and the single-band Otsu threshold are, on the median,
indistinguishable across the global pilot, and the dual-polarisation advantage
appears only at low A/P (Section~\ref{sec:res_fourway}). The fixed-training
dual SVM, by contrast, ranks below the single-band threshold, because it
cannot follow the scene-to-scene changes in backscatter that a per-scene
threshold absorbs by construction. The operational lesson is that the choice
worth making is not how many polarisations to use, but whether the decision
rule is re-estimated for each acquisition. This also settles the objection that
a fixed training year is arbitrary: retraining per scene removes the dependence
on any baseline year and, on the pilot, is at worst neutral and on the median a
small gain (Sections~\ref{sec:res_fourway} and~\ref{sec:res_sicily}).

A second consequence concerns cost, and here the picture is more nuanced than
the usual ``dual-pol is expensive'' assumption. Polarisation count itself is
cheap: a dual SVM trained once is cheaper than the single-band Otsu, whose
per-scene histogram dominates its cost. What is expensive is adaptivity.
Re-training the dual SVM on every scene, which is what makes it accurate, also
makes it the costliest method, somewhat above the VV Otsu
(Section~\ref{sec:res_cost}). Cost and accuracy therefore trade off, and because
the dual advantage is confined to low A/P, the two can be reconciled by using
the per-scene dual SVM only where it pays and the cheaper per-scene VV Otsu
elsewhere. This is the same A/P-based selection rule that the accuracy results
imply, now justified on cost as well.

\subsection{A/P as an a-priori screening tool}

Because A/P is computed from a water-body polygon that is already available in
global databases (GRanD, GRDL, or the JRC extent), it can be evaluated in
advance of processing. Its value is that it predicts the geometric ceiling
on accuracy: high-A/P reservoirs are tracked reliably, whereas low-A/P
reservoirs are only permitted, not guaranteed, to fail
(Section~\ref{sec:res_ap}). This supports two practical uses. First, screening:
a water manager can rank a portfolio of reservoirs by expected SAR reliability
and direct field effort or complementary altimetry to the low-A/P systems.
Second, classifier selection: because the dual-polarisation advantage
concentrates at low A/P, the extra polarisation is worth adding only for the
compact-shoreline minority, and a single-band threshold is adequate elsewhere.
Both uses turn a descriptive geometric index into an operational decision rule.
The screening is also inexpensive to deploy, since the training samples and the
reservoir boundary are drawn automatically from global water and land-cover
masks, which removes the manual delineation that would otherwise limit transfer
to new regions.

\subsection{Geometric and radiometric limits are distinct}

A/P bounds accuracy from geometry alone and is blind to radiometric failures,
and the pooled data make this explicit. The two high-A/P reservoirs with low
skill (Rockport and Boegoeberg) fail for reasons A/P cannot see, namely ice
cover and extreme drawdown (Section~\ref{sec:res_ap}). Wind, the failure mode
most often invoked for SAR water mapping, does not organise along geometry and
does not explain the between-method differences: the expected Bragg signature
is absent across the pilot, and A/P remains the only significant predictor of
the accuracy gap (Section~\ref{sec:res_wind}). We read this as evidence that
the geometric and radiometric axes are genuinely separate. A/P is a clean and
transferable predictor of the geometric limit, whereas the radiometric residual
is real but site- and date-specific and is not removed by the choice of
polarisation. Addressing it would require a different lever, such as a dynamic
threshold on the cross-polarised or ratio channel rather than on VV.

\subsection{Limitations}

Several limitations bound these conclusions. The study addresses surface area
only. Converting area to storage volume requires area--elevation--volume
information that for many reservoirs is outdated, and we leave that conversion,
together with the updating of those curves, to future work. The global tier is
referenced against the 30~m JRC product, which is an inter-product comparison
rather than ground truth; we mitigate this with the four Sicilian reservoirs
validated against 3~m PlanetScope near-truth, which reproduce the global
findings, but a broader near-truth validation would strengthen them. The A/P
classes used for display (low, medium, high) are a reading aid rather than a
formal taxonomy, and the underlying relationship is continuous. Finally, the
static A/P used for screening describes the full pool; at very low water the
effective shoreline complexity changes, which the dynamic A/P captures but the
a-priori index does not.

% PROVISIONAL: the global four-way comparison is currently N=28; the 10
% complementary reservoirs are being processed and will bring it to 38.